% Part: sets-functions-relations
% Chapter: relations
% Section: special-properties

\documentclass[../../../include/open-logic-section]{subfiles}

\begin{document}

\olfileid{sfr}{rel}{prp}
\olsection{د اړيکو خاص خاصيتونه}

\begin{intro}
د اړيکو ځينې ډولونه دومره عام دي چې خاص نومونه ورکړل شوي دي۔
د بېلګې په توګه، $\le$ او $\subseteq$ دواړه په خپلو خپلو ساحو کښې
(مثلاً، د $\le$ دپاره په $\Nat$ او د $\subseteq$ دپاره په $\Pow{A}$
کښې) په يو شان طريقو اړيکه جوړوي۔ د دې اړيکو د ورته والي او توپير
د دقيق معلومولو دپاره، د هغو خاصو خاصيتونو له مخې ئې وېشو چې
اړيکې ئې لرلے شي۔ څرګندېږي چې د دې ځينو خاصو خاصيتونو (ترکيبونه)
په تېره بيا مهم دي: ترتيبونه او د هم ارزښتۍ اړيکې۔
\end{intro}

\begin{defn}[انعکاس]
اړيکه $R \subseteq A^2$ هله او يوازې هله \emph{انعکاسي} ده چې
د هر $x \in A$ دپاره $Rxx$ وي۔
\end{defn}

\begin{defn}[انتقال]
اړيکه $R \subseteq A^2$ هله او يوازې هله \emph{انتقالي} ده چې
کله هم $Rxy$ او $Ryz$ وي، نو $Rxz$ هم وي۔
\end{defn}

\begin{defn}[تناظر]
اړيکه $R \subseteq A^2$ هله او يوازې هله \emph{متناظره} ده چې
کله هم $Rxy$ وي، نو $Ryx$ هم وي۔
\end{defn}

\begin{defn}[ضد تناظر]
اړيکه $R \subseteq A^2$ هله او يوازې هله \emph{ضد متناظره} ده چې
کله هم $Rxy$ او $Ryx$ دواړه وي، نو $x=y$ وي (په نورو ټکو: که
$x\neq y$ وي، نو يا $\lnot Rxy$ وي او يا $\lnot Ryx$)۔
\end{defn}

\begin{explain}
په متناظره اړيکه کښې $Rxy$ او $Ryx$ تل يوځاے وي، يا يو هم نۀ وي۔
په ضد متناظره اړيکه کښې $Rxy$ او $Ryx$ يوازې هله يوځاے کېدلے شي
چې $x = y$ وي۔ پام وکړئ چې دا خبره د $x = y$ په وخت د $Rxy$ او
$Ryx$ موجوديت \emph{نۀ لازموي}؛ يوازې دا چې مخه ئې نۀ نيول کېږي۔
نو ضد متناظره اړيکه انعکاسي کېدلے شي، خو هره ضد متناظره اړيکه
انعکاسي نۀ وي۔ دا هم په پام کښې وساتئ چې ضد متناظره کېدل او محض
متناظره نۀ کېدل بېل شرطونه دي۔ په رښتيا، يوه اړيکه په يوه وخت
متناظره او ضد متناظره دواړه کېدلے شي (مثلاً، د عينيت اړيکه داسې ده)۔
\end{explain}

\begin{defn}[تړلتيا]
اړيکه $R \subseteq A^2$ \emph{تړلې} ده که د ټولو $x,y\in A$
دپاره، که $x \neq y$ وي، نو يا $Rxy$ وي او يا $Ryx$۔
\end{defn}

\begin{prob}
د داسې اړيکو بېلګې ورکړئ چې (الف) انعکاسي او متناظرې وي، خو انتقالي
نۀ وي؛ (ب) انعکاسي او ضد متناظرې وي؛ (ج) ضد متناظرې او انتقالي وي،
خو انعکاسي نۀ وي؛ او (د) انعکاسي، متناظرې او انتقالي وي۔ پر عددونو
يا سټونو باندې اړيکې مۀ کاروئ۔
\end{prob}

\begin{defn}[د انعکاس نفي]
اړيکه $R \subseteq A^2$ \emph{نفي انعکاسي} بللے شي که د هر
$x \in A$ دپاره $Rxx$ نۀ وي۔
\end{defn}

\begin{defn}[نامتناظرتيا]
اړيکه $R \subseteq A^2$ \emph{نامتناظره} بللے شي که د هېڅ جوړې
$x,y\in A$ دپاره $Rxy$ او $Ryx$ دواړه ونۀ لرو۔
\end{defn}

پام وکړئ: که $A \neq \emptyset$ وي، نو په $A$ باندې هېڅ نفي انعکاسي
اړيکه انعکاسي نۀ ده، او په $A$ باندې هره نامتناظره اړيکه ضد متناظره
هم ده۔ خو داسې $R \subseteq A^2$ شته چې نۀ انعکاسي دي او نۀ نفي
انعکاسي، او داسې ضد متناظرې اړيکې شته چې نامتناظرې نۀ دي۔

\end{document}
